\chapter{Methodology and Approach}
\label{cha:MethodologyAndApproach}

\section{Dataset}
% things to mention:
% ids only from the ones of unarXive 2022
% randomly select paperId, get References and specter vectors from semantic scholar
% filter to references in unarXive 2022 dataset
% calculate cosine similarity between papers
% filter for cosine similarity > 0,7 and <0,9
% select 4 additional papers
% if not enough are valid, skip this question

% calculate adaptive individual character limit 
     % longer papers get more leftovers: adaptiveCharacterLimit = int(PAPER_CHARACTER_LIMIT + (len(text) / totalCharsInAbove50kPapers) * charactersLeftFromSub50kPapers)
% take half of individual character limits from beginning/ end of paper (highest information density)
% prompt LLM 
    % extract 10 topic overlaps, contradict or complementing, identify entities/concepts    
    % generate 1 question - answer pair 
    % requires information in multiple papers, synthesizing, resolving, applying , evaluating information
    % context independent / standalone answerable
    % complex scientific questions, not just simple/ definitial questions
    % question answer as precise as possible
    % add arxiv ids of used papers

% CoT prompt additional LLM for the question-answer pair and each individual paper text to determine whether it was used or not
% 

% talk about how the dataset looks
To curate a dataset of high-quality, expert-level scientific question-answer pairs, a fully automated pipeline is applied. 
It consists of three consecutive steps: document selection, question-answer generation and quality assurance in order to ensure that generated questions and answers are scientifically rigorous and strictly dependent on multi-document information synthesis in order to be answered, while also reliably providing associated ground-truth documents alongside them.
\subsection{Document selection}
To match the underlying document corpus, selectable papers are restricted to those of the unarXive 2022 dataset, yet for other applications, this could be expanded at will. A randomly selected anchor paper as well as its cited references serve as the foundation for generating a single question-answer pair. For each paper, a document-level SPECTER \cite{specter} vector embedding is retrieved via the Semantic Scholar academic graph API \cite{semanticScholarApi}. Leveraging citation graphs, these vectors provide dense representations of a paper's title and abstract, which are subsequently used to calculate the cosine similarity and thus the semantic similarity of each paper with the anchor document. To ensure topical alignment while excluding highly redundant work, a subset of candidate papers with a cosine similarity between 0.7 and 0.9 is selected. Four of these papers are then chosen at random to be used alongside the anchor paper for generation. Restricting the model input to five papers in total provides a multitude of topical overlaps for the model to choose from without resulting in exceedingly long prompts, mitigating the risk of "Lost in the Middle" effects. If a chosen anchor paper fails to provide at least four appropriate references, it is considered insufficient to be used in the dataset.
% rephrase realized, rephrase docuemnt usage verification
% justify 0.7 to 0.9
% maybe dont mention lost in the middle just here

\subsection{Question-Answer generation}
While modern LLMs offer increasingly large context windows, usually even above 100,000 tokens, empirical studies have shown severe performance degradation in tasks like complex reasoning, data retrieval and logic problems, accompanied by drastically increased hallucination \cite{contextIsWhatYouNeed}. In the generation step, the pipeline applies a moonshotai/Kimi-K2.6 \cite{moonshotKimi2.6} model, supporting a context window of approximately 256,000 tokens. In practice, however, the "effective context window", in which models only suffer from minimal performance degradations with increasingly long inputs, has been shown by empirical studies \cite{ruler} and practitioner reports to be drastically lower, with some studies estimating it as low as 40\% of the maximum supported context window \cite{intelligenceDegredation}. To ensure high reasoning capabilities, the pipeline applies an adaptive character limit for each paper text, enforcing a total upper bound of 200,000 characters for all papers combined, which conservatively translates to roughly 66,000 tokens for scientific paper texts \cite{howLongIsAPieceOfString}, or about 26\% of Kimi-K2-6's supported context size.
% meh 
Papers are assigned input space according to the following formula, allowing for full and efficient utilization of the context window, even if papers are highly variable in length:
\begin{equation}
C_{adaptive} = \left\lfloor C_{base} + \left( \frac{l_{paper}}{\sum_{p \in P_{above}} len(p)} \right) \times \sum_{s \in P_{below}} (C_{base} - len(s)) \right\rfloor
\end{equation}
Paper texts longer than their adaptive character limit are then head- and tail-truncated, retaining the first and last $\lfloor\frac{C_{adaptive}}{2}\rfloor$ characters, which are the most information-dense parts of scientific papers.

Subsequently, the LLM is prompted using a two-step strategy. Firstly, the model is tasked with selecting the best pair of papers (A, B), where paper A provides an intermediate entity, method, dataset or a similar element, and paper B evaluates, extends, contradicts or relies on A's intermediate element. Following that, the LLM is tasked with generating a single question-answer pair based on its selected pair, whose answer can only be derived by following a sequential reasoning process across both papers, applying paper A's intermediate element to paper B. Additionally, the model is instructed to answer the question as precisely as possible, providing a precision-biased subset of the underlying papers used during generation alongside.
% IMPORTANT FIX THE LIMIT IS 25000 TODO TODO TODO
% Lost in the Middle: How Language Models Use Long Contexts" (Liu et al., 2023).

% hallucinating used papers
% lazily attributing to many
% judge every single paper if really used, 
% CoT prompting, 
\subsection{Quality assurance}
\label{subsec:quality_assurance}
In order to ensure the dataset's integrity, two separate evaluation steps are applied, both using an openai/gpt-oss-120b \cite{gptoss120B} model.
The first judge evaluates each Question-Answer pair along six distinct dimensions, providing a judgement on how well each criterion is satisfied on a three-point ordinal scale. The complete prompt instructions for each, are provided in the appendix.
\begin{enumerate}
      \item  \textbf{Multi-Hop Validity} assesses whether answering the question requires integrating information from both papers. Thus, this criterion is only satisfied if neither paper provides sufficient information for deriving the complete answer.
      A question only satisfies this criterion if neither paper alone provides sufficient evidence to derive the complete answer (\ref{lst:MultihopValidityPrompt}).
      \item  \textbf{Dependency Strength} evaluates whether the steps in the reasoning chain are sequentially dependent in order for the conclusion to be drawn. Thus, the second reasoning step cannot be solved without first applying or interpreting the first one (\ref{lst:DependencyStrengthPrompt}).
      \item  \textbf{Non-Decomposability} verifies that questions cannot be broken down into independent single-hop sub-questions, allowing for reasoning shortcuts or simply combining sub-question answers. Instead, reasoning steps must be connected in such a way that separately solving them would not be sufficient for deriving an answer (\ref{lst:NonDecomposabilityPrompt}).
      \item  \textbf{Evidence Distribution} examines whether both papers contribute distinct and necessary evidence, ensuring that required evidence is spread across both papers and that questions cannot be answered by using a single dominant source with redundant information (\ref{lst:EvidenceDistributionPrompt}).
      \item  \textbf{Answerability} assesses whether the answer is fully supported and obtainable by only relying on information present in the papers, without requiring external knowledge or speculative inference not explicitly supported by either paper. 
      Is the answer fully supported by the provided evidence (\ref{lst:AnswerabilityPrompt})?
      \item  \textbf{Decontextualization} validates that the question is self-contained and unambiguous in an open-domain setting. Thus, questions cannot contain context-specific references and must instead introduce relevant concepts, methods or entities by themselves in order to satisfy this criterion (\ref{lst:decontextualizationPrompt}).
\end{enumerate}
% todo explain this more, better wording

An analysis of disconnected reasoning by Trivedi et al. \cite{isMultihopInDireCondition} demonstrated that many alleged multi-hop question-answering datasets do in fact not require true multi-hop reasoning, but can instead be solved through reasoning shortcuts, partial evidence or decomposition into independent sub-questions. Using a methodology inspired by the DIRE \cite{isMultihopInDireCondition} framework, the second judge verifies that questions cannot be answered under incomplete evidence, namely a subset $P' \subset P$ of ground truth papers $P$, but instead require all ground truth documents for a complete reasoning chain and answer. Thus, the test outcome satisfies: $$\left( \forall P' \subset P : \neg\text{Pass}(P') \right) \land \text{Pass}(P)$$

The final dataset is then obtained by filtering question-answer pairs to only those successfully passing both judges. In order to verify the pipeline and the judges, human verification will be applied on a subset.


\section{Cited text span extraction}
Constrained by the SQuAI system's attribution granularity, which natively provides exactly one source document per answer sentence, the source span extraction task in this work is formulated as a post-hoc attribution step for providing one source span per answer sentence. The retrieval methods themselves, however, are applicable to other systems with alternative attribution granularities.
This constitutes a limitation for this specific system, especially in cases where one attribution scope (e.g., a sentence) contains multiple claims requiring evidence, but also avoids additional ambiguity introduced by finer attribution granularity. For the purpose of comparing different extraction methods under identical conditions and constraints, however, this granularity is sufficient in order to measure whether extraction variants are capable of identifying relevant evidence, with results being directly comparable. 
% concern, does that really give good measurements, say that sentence level is already a good level of granularity and does not impact that much how good the methods perform - may not effect all sentences badly the same - some more some less

As previously discussed, cited text span attribution can be approached through a number of different paradigms, with methods being introduced and applied under vastly different circumstances. In the following, a hierarchy of increasingly complex post-hoc extraction methods that will be applied on a unified dataset is introduced. All methods except the native SQuAI span extraction and the generative LLM method receive tuples composed of an answer sentence and a set of eligible source spans $\mathcal{W}(P)$ from the paper $P$, annotated as the factual basis for that sentence during question answering. The set of source spans is composed of floating windows of contiguous sentences in the source document, up to a length of five, and can be denoted as: 

% \(D = (s_1, s_2, \ldots, s_n)\). Candidate spans are then constructed as
% floating windows of consecutive sentences with lengths from one to five:
\[
\mathcal{W}(P)
=
\bigcup_{k=1}^{5}
\left\{
s_i \oplus s_{i+1} \oplus \cdots \oplus s_{i+k-1}
\;\middle|\;
1 \leq i \leq n-k+1
\right\},
\]
where \(k\) specifies the window size and \(\oplus\) denotes string concatenation
with whitespace.

\begin{table}[H]
\centering
\resizebox{\textwidth}{!}{
\begin{tabular}{ccccc}
    \hline
    \textbf{\makecell{Extraction \\ Method}} & 
    \textbf{\makecell{Floating\\windows}} & 
    \textbf{\makecell{Retrieval\\Category}} & 
    \textbf{\makecell{Candidate Selection\\(Top 10)}}  & 
    \textbf{\makecell{Final Selection\\(Top 1)}} \\ \hline
    1  & $\times$  & Native        & -                                                               & SQuAI \\
    2  & \checkmark & Dense         & -                                                               & Bi-encoder \\
    3  & \checkmark & Sparse        & -                                                               & BM25 \\             
    4  & \checkmark & Hybrid        & Bi-encoder                                                      & BM25 \\ 
    5  & \checkmark & Hybrid        & BM25                                                            & Bi-encoder \\                                 
    6  & \checkmark & Hybrid        & -                                                               & BM25 \& Bi-encoder \tablefootnote{\ using Reciprocal Rank Fusion\label{fn:hybrid}}\\                                                             
    7  & \checkmark & Cross-encoder & Bi-encoder                                                      & Cross-encoder\\
    8  & \checkmark & Cross-encoder & BM25                                                            & Cross-encoder \\                                                             
    9  & \checkmark & Cross-encoder & BM25 \& Bi-encoder \textsuperscript{\ref{fn:hybrid}} & Cross-encoder\\                                                             
    10 & $\times$ & Generative    & - & LLM \\   \hline                                                    
\end{tabular}%
}
% todo check if squai really has no floating windows
% TODO better formattig here
    \caption{Cited text span extraction methods overview}
\end{table}

% todo talk about performance
% TODO
% assing a clear hypothesis what is tested to each of the methods
\subsection{Native SQuAI extraction baseline}
The first method relies on SQuAI's native text span extraction mechanism, utilizing simple lexical matching. Since it is only capable of providing one span per paper, even if it is cited repeatedly in different contexts, it is limited by design. Nevertheless, representing the attribution behaviour of the SQuAI system, it serves as a native baseline for further extraction pipelines to compare against.
\subsection{Dense retrieval}
Applying a bi-encoder, this method selects the best-fitting source span out of all floating windows based on the similarity of the span's dense vector representation with that of the answer sentence it is attributed to. It measures the suitability and sufficiency of relying exclusively on semantic similarity for source span attribution. As a dense retrieval baseline, it is capable of identifying appropriate source spans with different wording, but is prone to overlooking precise terminology. 
\subsection{Sparse retrieval}
The third method utilizes BM25, reranking source spans based on their lexical similarity with the corresponding answer sentence. The highest-scoring span is then selected as the final attributed span. Representing the sparse retrieval baseline, BM25 is especially effective with precise terminology, but less robust to paraphrasing. Its purpose is to test whether focusing on specific terms, entities or important keywords alone proves effective for cited text span extraction.

\subsection{Hybrid retrieval }
Hybrid methods aim to combine sparse and dense retrieval, leveraging their distinct characteristics. Thus, these variants evaluate whether the retrieval performance can be enhanced by combining both approaches into a single method. 
The fourth method applies a two-stage retrieval method, with the bi-encoder selecting the top ten candidate windows and BM25 selecting the final span, while method five reverses their order. While the former prioritizes semantic recall during candidate selection and lexical precision during final span selection, the latter reverses this order. Comparing both allows for analysis of whether either of those paradigms is better suited for candidate or final span selection. 
The sixth method combines both lexical and semantic matching into a single score using Reciprocal Rank Fusion (RRF). Instead of sequential application, this variant evaluates hybrid retrieval performance without an imposed order between the components themselves. Thus, this combination allows for assessment of whether the paradigms capture complementary aspects of the attribution task, resulting in performance enhancements when applied at the same time. 

\subsection{Cross-Encoder reranking}
Overcoming the limitation of BM25 and bi-encoder retrieval processing the answer sentence and source span in isolation, the following methods introduce cross-encoders, allowing for joint processing of both. This promises improvements, especially for context windows where fine-grained semantic differences such as negations or conditions as well as important entities or numerical values determine whether the answer sentence is merely related, contradicted or truly supported. These variants therefore evaluate whether joint answer-span comparison translates into better evidence selection in the form of cited text spans.
However, as cross-encoders are far more computationally expensive than previous methods and thus not feasible to run on the whole set of floating window candidate spans, these methods instead use different candidate selectors, with the cross-encoder selecting the best-fitting final span for the answer sentence.
Methods seven and eight apply semantic and lexical matching, respectively, while method nine utilizes both together as a candidate selector using RRF. Comparing these variations allows assessing which candidate selection strategy provides the best basis for a cross-encoder and therefore improves overall attribution quality.

\subsection{Generative LLM}
The last method uses a generative LLM to extract the best supporting span from the source document. Unlike previous methods, it does not rerank candidate spans, but instead prompts the model to identify the best-suited text span directly. Representing a generative approach, it is used to test whether an LLM's reasoning capabilities can improve evidence selection beyond retrieval and reranking, serving as a high-cost reference method.
% todo append prompt

\section{System evaluation}
The performance of the text span extraction methods can only be evaluated conditionally on the outputs of preceding pipeline steps, since post-hoc attribution is heavily dependent on their results. The evaluation is therefore split into assessing the suitability of the upstream inputs for text span extraction and evaluating the quality of the extracted spans themselves. This separation prevents text-span evaluation results from being confounded by errors introduced before span extraction.
% add performance evaluation?

\subsection{Upstream system evaluation}
% TODO 
% split error sources 
    % - the system retrieved the wrong paper,
    % - the system retrieved the right paper but the wrong paragraph,
    % - the answer sentence itself is factually wrong,
    % - the attribution method failed although the evidence was present,
    % - or the evaluation metric is too weak to recognize a good span.


Upstream inputs leading to bad span attributions can be flawed for a number of reasons, including:
% add method of extraction
\begin{enumerate}
    \item Retrieval failure: The system retrieves or filters the wrong or insufficient subset of papers.
    \item Generation failure: The system generates a factually wrong, overgeneralized, hallucinated or insufficiently supported answer.
    \item Document attribution failure: The system fails to attribute or incorrectly attributes source documents to answer sentences.
    % \item sentence contains multiple claims
\end{enumerate}

These must be addressed separately when evaluating the attributed spans. 

Retrieval failure can be analyzed by comparing the retrieved documents with those used during question generation, making it possible not only to evaluate the retrieval capabilities of the SQuAI system itself using precision, recall and F1-scores, but also the conditional performance of attribution methods given only unrelated, partial or complete evidence sources.

Generation failure is examined by comparing the system-generated answer with the ground-truth answer provided in the dataset. This comparison evaluates whether the system's answer covers at least the factual content of the reference, including all relevant facts and without introducing contradictory information. Exact lexical overlap is not required, as the comparison is performed by an LLM-as-a-judge, whose reliability is validated through human annotation on a small subset.

Document attribution failure can be evaluated by running extraction methods on all retrieved source documents, regardless of the sentence-specific document annotation by the answering LLM. This allows for comparing the best spans from both the annotated document and the whole set of retrieved papers. If no sufficient span can be extracted from the former, while the latter provides one, the failure indicates incorrect document attribution rather than a failure of the span extraction method itself. 

\subsection{Span evaluation}
% not much expierience with the judgment here yet, so preliminary
The attributed spans are evaluated across multiple distinct dimensions in span-sentence pairs, with the goal of assessing whether they provide sufficient evidence for claims made in their respective answer sentence. This evaluation applies both an LLM-as-a-judge approach with RAGAS-inspired metrics, such as faithfulness, answer relevance and context relevance \cite{ragas}, as well as an NLI model to check for logical entailment between the span and the answer sentence.
Additionally, the runtime required to extract spans will be compared across methods, allowing the evaluation to capture the trade-off between attribution quality and computational cost.



